\section{Implementation of Two-Finger Gripper Optimization}
\label{sec:two_finger_gripper_optimization}

This chapter presents the optimization framework developed for a two-finger gripper with a
spline-based gripping contour. The framework supports both evolutionary multi-objective optimization
and Bayesian multi-objective optimization. The evolutionary algorithms are implemented using the
\texttt{pymoo} library~\cite{pymoo}, while Bayesian optimization is implemented using the
\texttt{Ax} platform and its underlying \texttt{BoTorch} modeling
framework~\cite{olson2025ax,botorch}.

For both approaches, the evaluation of candidate gripper designs is delegated to an external service
through an API. The evaluation service constructs the gripper geometry, voxelizes the gripper and
workpieces, and calculates the grasp-quality objectives using a modified version of Snuggle-Pack\cite{snuggle_pack}.
This separation provides a modular architecture in which the optimization logic and the
domain-specific grasp evaluation remain decoupled.

The same decision variables, parameter bounds, geometry-generation pipeline, workpieces,
orientations, and grasp-quality objectives are used for both optimization approaches. This enables
their performance to be compared under a common evaluation budget and evaluation procedure.

\subsection{Decision Variables and Bounds for Two-Finger Gripper Problem}
\label{sec:two_finger_decision_variables}

A candidate gripper is represented by the decision vector
\begin{equation}
    \mathbf{x}
    =
    \begin{bmatrix}
        % d_{\mathrm{gn}} &
        l_{\mathrm{gn}} &
        h_1 &
        h_2 &
        h_3
    \end{bmatrix}^{\mathrm{T}},
\end{equation}
where the variables are transmitted to the evaluation API in the order
\begin{equation}
    \mathbf{x}
    =
    [
        % \texttt{depth\_gn},
        \texttt{length\_gn},
        \texttt{h\_p1\_spline},
        \texttt{h\_p2\_spline},
        \texttt{h\_p3\_spline}
    ]^{\mathrm{T}}.
\end{equation}

The decision variables and their bounds are summarized in Table~\ref{tab:two_finger_variables}.

\begin{table}[htbp]
\centering
\caption{Decision variables of the two-finger gripper optimization problem.}
\label{tab:two_finger_variables}
\begin{tabularx}{\textwidth}{@{}l l l >{\raggedright\arraybackslash}X@{}}
\toprule
\textbf{Variable} & \textbf{Configuration parameter} & \textbf{Lower bound} & \textbf{Upper bound}
\\
\midrule
% $d_{\mathrm{gn}}$ & \texttt{depth\_gn}       & 1 & 25 \\
$l_{\mathrm{gn}}$ & \texttt{length\_gn}      & 20 & 50 \\
$h_1$             & \texttt{h\_p1\_spline}   & 2 & 10 \\
$h_2$             & \texttt{h\_p2\_spline}   & 2 & 10 \\
$h_3$             & \texttt{h\_p3\_spline}   & 2 & 10 \\
\bottomrule
\end{tabularx}
\end{table}


\nomenclature{$\mathbf{x}$}{Decision vector of the two-finger gripper (see also $\mathbf{x}$ for generic decision variables)}
% \nomenclature{$d_{\mathrm{gn}}$}{Depth parameter in the common gripper interface}
\nomenclature{$l_{\mathrm{gn}}$}{Length of the gripping contour}
\nomenclature{$h_1,h_2,h_3$}{Spline contour heights}
All four decision variables are treated as continuous variables. Candidate values generated by
either optimization framework are restricted to the bounded search space
\begin{equation}
    \mathcal{X}
    =
    % [1,25]
    % \times
    [20,50]
    \times
    [2,10]^3.
\end{equation}
\nomenclature{$\mathcal{X}$}{Bounded decision space of the two-finger gripper (see also $\mathcal{X}$ for generic decision space)}

For evolutionary optimization, the variables are initially sampled within these bounds and
subsequently modified by the operators of the selected algorithm. For Bayesian optimization, Ax
initially evaluates the center of the search space and generates additional initialization
candidates using Sobol sampling. After sufficient observations have been collected, Ax transitions
to model-based candidate generation.

\subsection{Optimization Framework for Two-Finger Gripper Optimization}
\label{sec:two_finger_framework}

The two-finger gripper problem is solved using either an evolutionary multi-objective optimization
algorithm or Bayesian multi-objective optimization. Both methods use the same bounded decision
space, objective definitions, and external evaluation service.

The common optimization configuration contains:
\begin{itemize}
    \item Number of objectives $n_{\text{obj}}$
    \item Decision-variable definitions and bounds
    \item Optimization direction (minimization / maximization problem)
    \item Choice of optimization method
    \item Evaluation endpoint URL
\end{itemize}

For evolutionary optimization, the additional configuration parameters include:
\begin{itemize}
    \item Population size $N$
    \item Maximum number of generations $G_{\max}$
    \item Choice of evolutionary algorithm
\end{itemize}

For Bayesian optimization, the additional configuration parameters include:
\begin{itemize}
    \item Maximum number of candidate-generation batches $B_{\max}$, which is read from the same
          configuration field as the number of generations of the evolutionary algorithms
    \item Maximum number of candidates requested per batch $b$
\end{itemize}

In the Bayesian implementation, a batch represents a candidate-generation step rather than an
evolutionary generation. The maximum requested number of Bayesian evaluations is therefore
\begin{equation}
    E_{\max}^{\mathrm{BO}} = B_{\max}b.
\end{equation}
The actual number of evaluations may be smaller if the generation strategy returns fewer than $b$
candidates in a batch. \nomenclature{$B_{\max}$}{Maximum number of Bayesian candidate-generation batches (see also $B_{\max}$ in multi-suction optimization)}
\nomenclature{$b$}{Maximum number of Bayesian candidates requested per batch (see also $b$ in multi-suction optimization)}
\nomenclature{$E_{\max}^{\mathrm{BO}}$}{Maximum requested number of Bayesian objective evaluations (see also $E_{\max}^{\mathrm{BO}}$ in multi-suction optimization)}

\subsubsection{Supported Optimization Methods for Two-Finger Gripper Optimization}

The framework supports the evolutionary and Bayesian multi-objective optimization methods listed in
Table~\ref{tab:two_finger_algorithms}.
\begin{table}[H]
\centering
\caption{Supported multi-objective optimization methods for the two-finger gripper.}
\label{tab:two_finger_algorithms}

\small
\begin{tabularx}{\textwidth}{
    @{}
    p{0.20\linewidth}
    p{0.17\linewidth}
    >{\raggedright\arraybackslash}X
    @{}
}
\toprule
\textbf{Family} & \textbf{Identifier} & \textbf{Method} \\
\midrule
Evolutionary
& NSGA2
& Non-dominated Sorting Genetic Algorithm II \\

Evolutionary
& NSGA3
& Non-dominated Sorting Genetic Algorithm III \\

Evolutionary
& MOEAD
& Multi-Objective Evolutionary Algorithm based on Decomposition \\

Evolutionary
& CMOPSO
& Competitive Multi-Objective Particle Swarm Optimization \\

Evolutionary
& SMSEMOA
& S-Metric Selection Evolutionary Multi-Objective Algorithm \\

Bayesian
& AX-MBM
& Ax modular BoTorch model for multi-objective Bayesian optimization%
 \\
\bottomrule
\end{tabularx}
\end{table}

% For evolutionary algorithms requiring reference directions, such as NSGA-III and MOEA/D, uniform
% reference directions are generated using the method of Das and Dennis~\cite{das_dennis}:
% \begin{equation}
%     \Lambda =
%     \texttt{get\_reference\_directions}
%     (\text{``das-dennis''},\;n_{\text{obj}},\;n_{\text{partitions}}).
% \end{equation}
% \nomenclature{$\Lambda$}{Set of reference directions for decomposition-based algorithms (see also $\Lambda$ in multi-suction optimization)}
% \nomenclature{$n_{\text{obj}}$}{Number of objectives (see also $n_{\text{obj}}$ in multi-suction optimization)}

% The Bayesian implementation uses the Ax generation strategy identified at runtime as
% \texttt{Center\allowbreak+\allowbreak Sobol\allowbreak+\allowbreak MBM\allowbreak:\allowbreak fast}.
% This strategy consists of the following stages:

% \begin{enumerate}
%     \item Evaluation of the center of the bounded decision space
%     \item Sobol quasi-random sampling for initial space exploration
%     \item Model-based candidate generation using the Ax modular BoTorch model (MBM)
% \end{enumerate}

% The Center and Sobol stages do not use a surrogate model or an acquisition function. Once the
% generation strategy transitions to the MBM stage, Ax uses the available observations to fit
% probabilistic surrogate models and optimize an acquisition function for multi-objective candidate
% selection. For continuous multi-objective problems, this commonly involves Gaussian-process
% surrogates and a hypervolume-improvement acquisition function such as \texttt{qLogNEHVI}.

% The exact surrogate and acquisition-function classes are selected by Ax and can depend on the
% installed Ax and BoTorch versions, the number of observations, and the experiment configuration.
% Therefore, the active generation stage, surrogate model, and acquisition function should be recorded
% from the runtime configuration for reproducibility.

\subsection{Initialization of Two-Finger Gripper Optimization}
\label{sec:two_finger_initialization}

The initialization process depends on whether evolutionary or Bayesian optimization is selected.

\subsubsection{Evolutionary Initialization}

The selected evolutionary algorithm initializes a population of $N$ candidate grippers:
\begin{equation}
    \mathcal{U}^{(0)}
    =
    \left\{
        \mathbf{x}^{(0)}_1,
        \mathbf{x}^{(0)}_2,
        \ldots,
        \mathbf{x}^{(0)}_N
    \right\},
    \qquad
    \mathbf{x}^{(0)}_i\in\mathcal{X}.
\end{equation}
\nomenclature{$\mathcal{U}^{(0)}$}{Initial evolutionary population}
\nomenclature{$N$}{Evolutionary population size (see also $N$ in multi-suction optimization)}

Each candidate is sampled within the variable bounds. The complete population is evaluated through
the external API before the selected evolutionary algorithm performs its first selection and
variation step.

\subsubsection{Bayesian Initialization}

Bayesian optimization does not maintain an evolutionary population. Instead, it maintains a dataset
containing the completed gripper evaluations:
\begin{equation}
    \mathcal{D}_t
    =
    \left\{
        \left(
            \mathbf{x}_i,
            \mathbf{k}(\mathbf{x}_i)
        \right)
    \right\}_{i=1}^{t},
\end{equation}
where
\begin{equation}
    \mathbf{k}(\mathbf{x})
    =
    \begin{bmatrix}
        k_1(\mathbf{x}) &
        k_2(\mathbf{x}) &
        \cdots &
        k_{n_{\text{obj}}}(\mathbf{x})
    \end{bmatrix}^{\mathrm{T}}
\end{equation}
contains the grasp-quality KPIs returned by the evaluation service.
\nomenclature{$\mathcal{D}_t$}{Dataset of completed Bayesian trials after $t$ evaluations (see also $\mathcal{D}_t$ in multi-suction optimization)}
\nomenclature{$\mathbf{k}(\mathbf{x})$}{Natural grasp-quality objective vector (see also $\mathbf{g}(\mathbf{q})$ in multi-suction optimization for similar concept)}

The Ax generation strategy initially evaluates the center of the bounded search space. Additional
initialization candidates are generated using Sobol quasi-random sampling. Sobol sampling provides
broad and approximately space-filling coverage of the four-dimensional search space before a
surrogate model is introduced.

After the initialization requirements of the generation strategy have been satisfied, Ax transitions
to the model-based modular BoTorch model (MBM) stage. From this point onward, new gripper designs are selected using
surrogate predictions and predictive uncertainty.

\subsection{Objective Function Evaluation of Two-Finger Gripper via API}
\label{sec:two_finger_api_evaluation}

For every candidate, the decision variables are sent to the evaluation API as an ordered list. The
API inserts these values into a complete gripper configuration, from which the jaw geometry is
generated:
\begin{equation}
    \mathbf{x}
    \longrightarrow
    \mathbf{c}_{\mathrm{gripper}}(\mathbf{x})
    \longrightarrow
    \mathcal{G}(\mathbf{x}),
\end{equation}
where $\mathbf{c}_{\mathrm{gripper}}$ denotes the gripper configuration and
$\mathcal{G}(\mathbf{x})$ the resulting geometry.

Besides the decision variables, the configuration contains fixed parameters, including the contour
type, the basic jaw dimensions (depth, width, and length), the jaw angle, the stroke, the axial
offset of the contour, and the tool-center-point translation.
\subsubsection{Parametrization of the Jaw Contour}
\label{sec:jaw_contour_parametrization}

The contact contour of the rigid jaw is described by an interpolating spline in the contour plane,
which is spanned by the jaw depth direction and the jaw length direction. The decision vector
\begin{equation}
    \mathbf{x}
    =
    \begin{bmatrix}
        \ell_{\mathrm{GN}} &
        h_1 &
        h_2 &
        h_3
    \end{bmatrix}^{\top}
    \in \mathcal{X} \subset \mathbb{R}^{4}
\end{equation}
consists of the contour length $\ell_{\mathrm{GN}}$ and three depth parameters $h_1$, $h_2$, and
$h_3$. These four values define the five control points of the spline,
\begin{equation}
    \mathbf{p}_0 = (0,\,0),
    \qquad
    \mathbf{p}_j
    =
    \left(
        h_j,\;
        \tfrac{j}{4}\,\ell_{\mathrm{GN}}
    \right),
    \quad j = 1,2,3,
    \qquad
    \mathbf{p}_4 = (0,\,\ell_{\mathrm{GN}}),
\end{equation}
so that the control points are distributed at the relative contour positions
$0,\ 0.25,\ 0.50,\ 0.75$, and $1.0$ along the contour length. The first and last control point are
pinned to zero depth, which guarantees that the contour begins and ends flush with the planar jaw
face and that the resulting profile is a closed recess rather than an open step. The intermediate
depths $h_1$, $h_2$, and $h_3$ control how far the contour intrudes into the jaw at the quarter,
half, and three-quarter positions and therefore determine the shape of the form fit. The contour
curve is the interpolating spline
\begin{equation}
    \mathcal{S}(\mathbf{x})
    =
    \operatorname{spline}
    \left(
        \mathbf{p}_0,\ldots,\mathbf{p}_4
    \right),
\end{equation}
which passes through all five control points.

The spline is closed into a planar profile by straight segments up to a fixed cutting depth and is
then extruded through the full jaw depth to remove the enclosed material from the basic jaw body.
The geometry is constructed with CadQuery~\cite{cadquery2026} and exported to the mesh
representation required for evaluation. Since only $\ell_{\mathrm{GN}}$ and the three depth
parameters are varied, every candidate geometry differs from the others exclusively in the contact
contour, while all remaining jaw dimensions remain identical.

Each candidate geometry $\mathcal{G}(\mathbf{x})$ is generated and voxelized only once per
evaluation request and is subsequently reused for all configured grasps, so that the cost of the
geometry generation is independent of the number of objectives.

\nomenclature{$\mathbf{x}$}{Decision vector of the parallel-jaw contour}
\nomenclature{$\ell_{\mathrm{GN}}$}{Contour length of the jaw form fit}
\nomenclature{$h_j$}{Spline depth parameter at relative contour position $j/4$}
\nomenclature{$\mathbf{p}_j$}{Spline control point $j$ of the jaw contour}
\nomenclature{$\mathcal{S}(\mathbf{x})$}{Interpolating spline describing the jaw contact contour}
\nomenclature{$\mathcal{G}(\mathbf{x})$}{Gripper jaw geometry generated from $\mathbf{x}$}
\nomenclature{$\mathbf{c}_{\mathrm{gripper}}$}{Complete gripper configuration}

\subsubsection{Grasp Set Definition}
\label{sec:grasp_set}

A grasp specifies the pose in which the workpiece is presented to the jaws. In the evaluation
pipeline, the gripper remains stationary in its nominal pose, and the workpiece is reoriented for
each grasp.

Only the orientation component of the generated grasp poses is used, since the relative
gripper--workpiece position is determined by the collision- and proximity-based placement search
described in Section~\ref{sec:grasp_eval}. A grasp $g_i$ is therefore fully described by a spatial
rotation $\mathbf{R}_i \in \mathrm{SO}(3)$ applied to the workpiece,
\begin{equation}
    \mathcal{W}_i = \mathbf{R}_i\,\mathcal{W},
    \qquad
    i = 1,\ldots,n_g ,
\end{equation}
where $\mathcal{W}$ denotes the workpiece in its reference orientation. The translation components
$\mathbf{t}_i$ of the generated poses are discarded. The grasp set is

\begin{equation}
    \mathcal{R}
    =
    \left\{
        \mathbf{R}_1,\ldots,\mathbf{R}_{n_g}
    \right\},
    \qquad
    n_g = 10.
\end{equation}


Because the retained poses originate from the feasibility filtering and the farthest-point
clustering of the generation procedure, the orientations are neither restricted to a single rotation
axis nor obtained as independent samples of a parametric distribution. They are instead a
spatially spread subset of the physically admissible grasps of the given workpiece, which ensures
that the evaluated orientations are mutually dissimilar and at the same time realizable on the part
geometry.

The grasp set is generated once before the optimization starts and is stored as part of the static
configuration. It is therefore identical for all candidates and for all repetitions of a run, which
makes the objective values of different candidates directly comparable. The orientations are fixed
quantities and are not optimization variables.

Each optimization run is defined for exactly one workpiece together with its corresponding
predefined grasp set. The workpiece geometry and the grasp set $\mathcal{P}$ are constant throughout
a run, and all objectives of a run refer to the same workpiece evaluated in its $n_g$ predefined
orientations. Different workpieces are treated as separate optimization problems and are optimized
in separate runs, since each workpiece requires its own task-specific jaw contour. In the present
study, $n_g = 10$ orientations are configured.

\nomenclature{$\mathcal{W}$}{Workpiece geometry in its reference orientation}
\nomenclature{$\mathcal{W}_i$}{Workpiece geometry in the orientation of grasp $g_i$}
\nomenclature{$\mathbf{R}_i$}{Spatial rotation defining grasp $g_i$ (see also $\mathbf{R}_i$ in multi-suction grasp generation)}
\nomenclature{$\mathcal{R}$}{Set of predefined workpiece orientations of a two-finger optimization run}
\nomenclature{$n_g$}{Number of predefined grasps of a two-finger optimization run}

\subsubsection{Grasp Evaluation}
\label{sec:grasp_eval}

The generated gripper is evaluated separately for every grasp $g_m \in \mathcal{P}$ of the single
workpiece configured for the run. For each grasp, the gripper geometry and the correspondingly
oriented workpiece geometry are voxelized to obtain three-dimensional occupancy grids.

The grasp evaluation uses a modified version of Snuggle-Pack, a three-dimensional packing algorithm
that combines FFT-based spatial analysis with proximity heuristics~\cite{snuggle_pack}. In the
original method, correlations between voxel grids are used to identify collision-free object
placements. Here, the same principle is adapted to evaluate relative gripper--workpiece placements.

Let $s_{\mathcal{G}}$ and $s_{\mathcal{W}_m}$ denote the occupancy grids of the gripper and of the
workpiece in the orientation of grasp $g_m$. Their spatial correlation produces a collision map:
\begin{equation}
    \zeta_{\mathcal{G},\mathcal{W}_m}(\boldsymbol{\rho})
    =
    \operatorname{DFT}^{-1}
    \left(
        \overline{\operatorname{DFT}(s_{\mathcal{G}})}
        \,\operatorname{DFT}(s_{\mathcal{W}_m})
    \right)(\boldsymbol{\rho}).
\end{equation}
where $\boldsymbol{\rho}$ is a relative displacement and $\mathcal{F}$ denotes the Fourier transform.
Values greater than zero indicate overlapping voxels, whereas zero values identify collision-free
relative placements.

The collision map is combined with Gaussian proximity distributions to determine how closely the
fingers surround the workpiece at the feasible placements. These calculations are performed using
FFT-based convolution and GPU-accelerated PyTorch operations. The resulting grasp quality of grasp
$g_m$ is represented by the KPI \texttt{max\_prox}, denoted $k_m(\mathbf{x})$. If the opening
required for grasp $g_m$ cannot be achieved within the available gripper stroke, the corresponding
KPI is set to zero. The stroke feasibility is thus assessed per grasp, so that an infeasible
orientation penalizes only its own objective and leaves the remaining objectives unaffected.

\nomenclature{$s_{\mathcal{G}}$}{Gripper voxel occupancy grid}
\nomenclature{$s_{\mathcal{W}_m}$}{Voxel occupancy grid of the workpiece in grasp orientation $m$}
\nomenclature{$\zeta_{\mathcal{G},\mathcal{W}_m}$}{Gripper--workpiece collision map for grasp $g_m$}
\nomenclature{$\boldsymbol{\rho}$}{Relative displacement between the voxel grids}
\nomenclature{$\operatorname{DFT}$}{Discrete Fourier transform}

\subsubsection{Objective Evaluation}
\label{sec:two_finger_objective_evaluation}

The evaluation service returns one grasp-quality KPI per configured grasp. Because a single
workpiece and its associated grasp set are used per optimization run, the number of objectives is
determined directly by the size of that grasp set:
\begin{equation}
    n_{\text{obj}} = n_g .
\end{equation}
The individual KPIs are not aggregated into a scalar
value; each grasp retains its own objective, which preserves the information about which
orientations a given contour favours.

Since a larger proximity KPI represents a better grasp result, the natural optimization problem is
\begin{equation}
    \max_{\mathbf{x}\in\mathcal{X}} k_m(\mathbf{x}),
    \qquad
    m=1,\ldots,n_{\text{obj}},
\end{equation}
with the natural objective vector
\begin{equation}
    \mathbf{k}(\mathbf{x})
    =
    \begin{bmatrix}
        k_1(\mathbf{x}) &
        k_2(\mathbf{x}) &
        \cdots &
        k_{n_{\text{obj}}}(\mathbf{x})
    \end{bmatrix}^{\mathrm{T}}
    \in \mathbb{R}^{n_g}.
\end{equation}
Because the four decision variables define a single jaw contour that is shared by all grasps, the
objectives are generally conflicting: a contour whose depth profile conforms closely to the
workpiece in one orientation does not necessarily conform in the remaining ones. The optimization
therefore yields Pareto-optimal contours that represent different compromises in geometric
conformity across the ten predefined orientations.

\nomenclature{$k_m(\mathbf{x})$}{Grasp-quality KPI (\texttt{max\_prox}) of grasp $g_m$}
\nomenclature{$\mathbf{k}(\mathbf{x})$}{Natural maximization objective vector}
\nomenclature{$n_{\text{obj}}$}{Number of objectives}

The objective-direction convention is handled differently by the two optimization frameworks. Since
\texttt{pymoo} assumes minimization, the evolutionary algorithms negate the objective vector:
\begin{equation}
    \mathbf{f}_{\mathrm{EA}}(\mathbf{x})
    =
    -\mathbf{k}(\mathbf{x})
    =
    \begin{bmatrix}
        -k_1(\mathbf{x}) &
        -k_2(\mathbf{x}) &
        \cdots &
        -k_{n_{\text{obj}}}(\mathbf{x})
    \end{bmatrix}^{\mathrm{T}}.
\end{equation}

In the Bayesian implementation, the values returned by the external evaluation service are used
unchanged, since the maximization direction is declared in the \texttt{Ax} optimization
configuration:
\begin{equation}
    \mathbf{f}_{\mathrm{BO}}(\mathbf{x})
    =
    \mathbf{k}(\mathbf{x}).
\end{equation}
Consequently, the observed KPI values do not need to be negated before they are attached to the
corresponding Ax trial, which preserves the physical meaning of the KPI values when the final
Pareto frontier is returned.

\nomenclature{$\mathbf{f}_{\mathrm{EA}}(\mathbf{x})$}{Minimization-form objective vector supplied to the evolutionary algorithms}
\nomenclature{$\mathbf{f}_{\mathrm{BO}}(\mathbf{x})$}{Objective observations supplied to Ax}

For performance comparisons, the results of both optimization frameworks are converted to the common
natural maximization convention $\mathbf{k}(\mathbf{x})$, so that all reported objective values
correspond to the ten grasp-quality KPIs in their original direction and magnitude.

\subsection{Optimization Loop for Two-Finger Gripper}
\label{sec:two_finger_optimization_loop}

The external evaluation pipeline is shared by the evolutionary and Bayesian optimization approaches.
The methods differ primarily in how they use the completed evaluations to generate subsequent
candidates.

For every candidate, the common evaluation steps are:
\begin{enumerate}
    \item Send the candidate decision vector to the external evaluation API.
    \item Map the decision variables to the complete gripper configuration.
    \item Generate the spline-based gripper geometry using CadQuery.
    \item Voxelize the gripper and configured workpieces.
    \item Calculate FFT-based collision maps using the modified Snuggle-Pack procedure.
    \item Calculate one proximity KPI for each predefined orientation of the configured workpiece.
    \item Return the KPI vector to the selected optimization framework.
\end{enumerate}

\subsubsection{Evolutionary Optimization Loop}

In every evolutionary generation, the selected algorithm generates a population of candidate
decision vectors. The candidates are evaluated through the common evaluation pipeline, and the
resulting KPI values are negated before being supplied to \texttt{pymoo}.

The population update can be expressed abstractly as
\begin{equation}
    \mathcal{U}^{(g+1)}
    =
    \mathcal{A}
    \left(
        \mathcal{U}^{(g)},
        \mathbf{F}_{\mathrm{EA}}^{(g)}
    \right).
\end{equation}
where $\mathcal{A}$ represents the selected evolutionary algorithm and
$\mathbf{F}_{\mathrm{EA}}^{(g)}$ contains the minimization-form objective vectors evaluated in
generation $g$. \nomenclature{$\mathcal{A}$}{Selected multi-objective evolutionary algorithm (see also $\mathcal{A}$ in multi-suction optimization)}
\nomenclature{$g$}{Evolutionary generation index (see also $t$ in multi-suction optimization)}
\nomenclature{$\mathbf{F}_{\mathrm{EA}}^{(g)}$}{Evolutionary objective values in generation $g$ (see also $\mathbf{F}_{\mathrm{EA}}^{(g)}$ for evolutionary objective values)}

Depending on the selected algorithm, the next population is produced using operations such as
selection, crossover, mutation, decomposition, particle movement, or survival based on hypervolume
contribution.

\subsubsection{Bayesian Optimization Loop}

During the model-based stage, Ax and BoTorch use the completed trial dataset $\mathcal{D}_t$ to fit
probabilistic surrogate models:
\begin{equation}
    \widehat{\mathbf{k}}_t
    =
    \operatorname{Fit}_{\mathrm{BO}}(\mathcal{D}_t).
\end{equation}
where $\operatorname{Fit}_{\mathrm{BO}}$ denotes the surrogate-modeling procedure. 
\nomenclature{$\operatorname{Fit}_{\mathrm{BO}}$}{Bayesian surrogate-model fitting procedure}
\nomenclature{$\widehat{\mathbf{k}}_t$}{Probabilistic surrogate approximation after $t$ evaluations (see also $\widehat{\mathbf{g}}_t$ in multi-suction optimization)}

The surrogate models approximate the relationship between the gripper parameters and the
workpiece-specific KPI values. In addition to predicted objective values, they estimate predictive
uncertainty in regions containing few or no observations.

An acquisition function $\alpha(\mathbf{x}\mid\mathcal{D}_t)$ uses these predictions to balance:
\begin{itemize}
    \item Exploitation of gripper designs expected to improve the current Pareto frontier
    \item Exploration of uncertain regions of the bounded decision space
\end{itemize}

For a requested batch size $b$, the next candidate batch is generated as
\begin{equation}
    \mathcal{Q}_{t+1}
    =
    \operatorname*{arg\,max}_{
        \substack{
            \mathcal{B}\subset\mathcal{X}\\
            |\mathcal{B}|\leq b
        }
    }
    a_{\mathrm{BO}}(\mathcal{B}\mid\mathcal{D}_t).
\end{equation}
\nomenclature{$\mathcal{Q}_{t+1}$}{Next batch of two-finger gripper candidates}
\nomenclature{$\alpha$}{Bayesian acquisition function (see also $\alpha$ in multi-suction optimization)}

After all candidates in the generated batch have been evaluated, the dataset is updated as
\begin{equation}
    \mathcal{D}_{t+|\mathcal{Q}_{t+1}|}
    =
    \mathcal{D}_t
    \cup
    \left\{
        \left(
            \mathbf{x},
            \mathbf{k}(\mathbf{x})
        \right)
        \;\middle|\;
        \mathbf{x}\in\mathcal{Q}_{t+1}
    \right\}.
\end{equation}

After new trials have been completed, the surrogate models are updated using the expanded dataset,
and another candidate batch is generated. This process continues until the Bayesian evaluation
budget is exhausted.

Although multiple candidates can be requested in one Bayesian batch, the current implementation
evaluates the candidates one after another through the external API. The implementation therefore
performs \emph{sequential batch evaluation}, rather than parallel batch evaluation.
\subsection{Termination Criteria for Two-Finger Gripper Optimization}
\label{sec:two_finger_termination}

The termination criteria depend on the selected optimization approach.

\begin{table}[htbp]
\centering
\caption{Termination criteria for evolutionary and Bayesian optimization.}
\label{tab:two_finger_termination}
\begin{tabularx}{\textwidth}{@{}l l l >{\raggedright\arraybackslash}X@{}}
\toprule
\textbf{Approach} & \textbf{Criterion} & \textbf{Symbol} & \textbf{Description} \\
\midrule
Evolutionary & Design-space tolerance & $x_{\text{tol}}$ & Change in decision variables below
threshold \\

Evolutionary & Objective-space tolerance & $f_{\text{tol}}$ & Relative objective improvement below
threshold \\

Evolutionary & Maximum generations & $G_{\max}$ & Maximum number of generations reached \\

Evolutionary & Maximum evaluations & $E_{\max}$ & Maximum number of evaluations reached \\

Evolutionary & Hypervolume & $HV_{\max}$ & Hypervolume progression stagnates  \\

Bayesian & Maximum batches & $B_{\max}$ & Maximum number of candidate batches reached \\

Bayesian & Maximum evaluations & $E_{\max}^{\mathrm{BO}}$ & Bayesian evaluation budget exhausted \\
\bottomrule
\end{tabularx}
\end{table}

For evolutionary optimization, the design-space and objective-space tolerances are evaluated over a
sliding window of $w$ generations. \nomenclature{$x_{\text{tol}}$}{Design-space convergence tolerance (see also $x_{\text{tol}}$ in multi-suction optimization)}
\nomenclature{$f_{\text{tol}}$}{Objective-space convergence tolerance (see also $f_{\text{tol}}$ in multi-suction optimization)}
\nomenclature{$G_{\max}$}{Maximum number of evolutionary generations (see also $G_{\max}$ in multi-suction optimization)}
\nomenclature{$E_{\max}$}{Maximum number of evolutionary evaluations (see also $E_{\max}$ in multi-suction optimization)}
\nomenclature{$w$}{Number of generations in the convergence window (see also $w$ in multi-suction optimization)}

The current Bayesian implementation uses budget-based termination. It performs at most $B_{\max}$
candidate-generation batches and requests up to $b$ candidates in each batch. Therefore,
\begin{equation}
    E_{\max}^{\mathrm{BO}}=B_{\max}b.
\end{equation}



\subsection{Performance Monitoring for Two-Finger Gripper}
\label{sec:two_finger_monitoring}

\todo{ADD IGD+ information}
The hypervolume indicator is used to monitor both evolutionary and Bayesian
optimization. Before calculating hypervolume, all objective vectors are represented using the same
natural maximization convention:
\begin{equation}
    \mathbf{k}(\mathbf{x})
    =
    [k_1(\mathbf{x}),\ldots,k_{n_{\text{obj}}}(\mathbf{x})]^{\mathrm{T}}.
\end{equation}

For evolutionary optimization, a callback records the optimization state after each generation. Let
$\mathcal{K}_{\mathrm{EA}}^{(g)}$ denote the non-dominated KPI vectors available after generation
$g$. The corresponding hypervolume is
\begin{equation}
    \mathrm{HV}_{\mathrm{EA}}^{(g)}
    =
    \mathrm{HV}
    \left(
        \mathcal{K}_{\mathrm{EA}}^{(g)},
        \mathbf{r}
    \right).
\end{equation}

For Bayesian optimization, hypervolume can be calculated after every completed trial or after every
completed candidate batch. Let $\mathcal{K}_{\mathrm{BO}}^{(t)}$ denote the non-dominated KPI
vectors obtained from the first $t$ completed Bayesian trials. The corresponding hypervolume is
\begin{equation}
    \mathrm{HV}_{\mathrm{BO}}^{(t)}
    =
    \mathrm{HV}
    \left(
        \mathcal{K}_{\mathrm{BO}}^{(t)},
        \mathbf{r}
    \right).
\end{equation}
\nomenclature{$\mathcal{K}_{\mathrm{EA}}^{(g)}$}{Non-dominated KPI set after evolutionary generation $g$ (see also $\mathcal{F}^{(t)}$ in multi-suction optimization for similar concept)}
\nomenclature{$\mathcal{K}_{\mathrm{BO}}^{(t)}$}{Non-dominated KPI set after $t$ Bayesian evaluations (see also $\mathcal{F}^{(t)}$ in multi-suction optimization for similar concept)}
\nomenclature{$\mathrm{HV}_{\mathrm{EA}}^{(g)}$}{Evolutionary hypervolume after generation $g$ (see also $\text{HV}^{(t)}$ in multi-suction optimization)}
\nomenclature{$\mathrm{HV}_{\mathrm{BO}}^{(t)}$}{Bayesian hypervolume after $t$ completed evaluations (see also $\text{HV}^{(t)}$ in multi-suction optimization)}
\nomenclature{$\mathbf{r}$}{Fixed hypervolume reference point (see also $\mathbf{r}$ for reference point in theoretical background)}

Under the maximization convention, the fixed reference point $\mathbf{r}$ must be worse than the
objective values of interest. The same reference point and objective scaling must be used for all
compared methods. An increasing hypervolume generally indicates improvement in convergence toward
the Pareto front, coverage of the objective space, or both.

Evolutionary generations and Bayesian batches can contain different numbers of candidate
evaluations. Therefore, the primary comparison is performed using hypervolume as a function of the
cumulative number of objective evaluations:
\todo{Check this}
\begin{equation}
    \mathrm{HV}=\mathrm{HV}(E),
\end{equation}
where $E$ is the number of completed external API evaluations.
\nomenclature{$E$}{Cumulative number of completed external evaluations (see also $E_{\max}^{\mathrm{BO}}$ for Bayesian evaluation budget)}

The framework additionally records total wall-clock time because every function evaluation includes
API communication, CadQuery geometry generation, voxelization, and GPU-based FFT spatial analysis.



\subsection{Optional Hyperparameter Optimization for Two-Finger Gripper Problem}
\label{sec:two_finger_hyperparameter_optimization}

For the evolutionary algorithms, the framework optionally supports automated tuning of algorithm
hyperparameters, \todo{Add more info about which Hyperparamts}such as crossover probability and mutation rate, using a meta-optimization
approach. A \texttt{HyperparameterProblem} is constructed around the selected algorithm and
evaluated over multiple random seeds. The negative mean hypervolume across the runs is used as the
meta-objective:
\begin{equation}
    \min_{\boldsymbol{\theta}}
    \left\{
        -\frac{1}{|\mathcal{Z}|}
        \sum_{s\in\mathcal{Z}}
        \mathrm{HV}(\mathcal{K}_s,\mathbf{r})
    \right\}.
\end{equation}
\nomenclature{$\boldsymbol{\theta}$}{Evolutionary algorithm hyperparameters (see also $\boldsymbol{\theta}$ for hyperparameter optimization)}
\nomenclature{$\mathcal{Z}$}{Set of random seeds used for meta-optimization (see also $\mathcal{Z}$ for random seeds)}

The best hyperparameter configuration is subsequently applied to the main evolutionary optimization
run.

The current Bayesian implementation uses the generation strategy, surrogate configuration, and
acquisition-function configuration selected by Ax. A separate meta-optimization of Bayesian model
hyperparameters is therefore not currently performed. The user specifies the candidate batch size
and optimization budget, while Ax manages the transition from initialization to model-based
optimization.





\subsection{Output of Two-Finger Gripper Optimization}
\label{sec:two_finger_output}

The optimization produces a set of non-dominated gripper configurations. The Pareto-optimal set is
defined as
\begin{equation}
    \mathcal{P}^{*}
    =
    \left\{
        \mathbf{x}\in\mathcal{P}
        \;\middle|\;
        \nexists\,\mathbf{x}'\in\mathcal{P}:
        \mathbf{x}' \prec \mathbf{x}
    \right\},
\end{equation}
where $\mathbf{x}'\prec\mathbf{x}$ means that
\begin{equation}
    f_m(\mathbf{x}')\leq f_m(\mathbf{x})
    \quad
    \forall m
\end{equation}
and
\begin{equation}
    f_j(\mathbf{x}')<f_j(\mathbf{x})
    \quad
    \text{for at least one objective }j.
\end{equation}

The final output includes:
\begin{itemize}
    \item The Pareto-optimal decision vectors
    $\mathbf{x}=[l_{\mathrm{gn}},h_1,h_2,h_3]^{\mathrm{T}}$
    \item The corresponding objective KPI values
    \item The hypervolume progression
    $\{\mathrm{HV}^{(0)},\mathrm{HV}^{(1)},\ldots,\mathrm{HV}^{(T)}\}$
    \item Total optimization wall-clock time
\end{itemize}

The non-dominated solutions represent the best trade-offs between grasp quality across
the predefined orientations of the configured workpiece.

\nomenclature{$\mathcal{P}^{*}$}{Set of Pareto-optimal gripper configurations (see also $\mathcal{P}^*$ for Pareto-optimal solutions)}